\documentclass[10pt]{article}
\usepackage[letterpaper,margin=1in]{geometry}
\usepackage[T1]{fontenc}
\usepackage{lmodern,booktabs,microtype}
\usepackage[hidelinks]{hyperref}
\usepackage{enumitem}
\setlist{nosep,leftmargin=1.2em}
\setlength{\parindent}{0pt}
\setlength{\parskip}{3pt}
\pagestyle{empty}
\newcommand{\sect}[1]{\par\vspace{5pt}{\bfseries #1}\par\vspace{1pt}}
\begin{document}

{\large\bfseries ProtocolQA: Haiku 5.5 vs.\ GPT-6 Luna vs.\ GLM 5.3 Flash}\par
\vspace{2pt}

\sect{Eval design}
\textbf{Data.} The 108 public multiple-choice questions of LAB-Bench ProtocolQA (Laurent
et al., 2024; 27 more are held out). The authors and contracted biologists introduced
errors into published protocols by modifying or omitting steps. Each question gives the
full protocol (103--2,642 words) and an observed failure (e.g.\ low PBMC yield) and asks
which change fixes it. Options: one correct answer, 3--6 distractors, and an
``insufficient information'' abstain option (chance $\approx$ 19\%). The 20 open-answer
items were skipped (scoring needs an LLM judge; out of scope).

\textbf{What it measures.} Text-only protocol troubleshooting: tracing a symptom to the
faulty step and choosing the fix from fixed options, which takes domain knowledge and
careful reading of a long procedure. It does not test protocol design, open-ended
troubleshooting, tool or literature use, or real lab outcomes.

\textbf{Protocol.} Zero-shot with the original LAB-Bench chain-of-thought prompt, via
OpenRouter pinned to first-party endpoints. Reasoning effort \texttt{high}, 32k-token
output limit. 5 repeats per question with a seeded choice order shared by all models (540
calls per model). Coverage = share of calls answered (not abstained, blocked, or refused);
precision = correct/answered. CIs bootstrap over questions; pairs: sign-flip test.

\sect{Eval results}
\vspace{-4pt}\begin{center}\small
\begin{tabular}{lcccccc}
\toprule
Model & Accuracy [95\% CI] & Precision & Coverage & Blocked & Refused & Cost \\
\midrule
GLM 5.3 Flash & \textbf{0.678} [0.604, 0.750] & 0.696 & 0.974 & 0 & 0 & \$0.22 \\
Haiku 5.5     & 0.644 [0.565, 0.722] & 0.753 & 0.856 & 5 & 0 & \$0.38 \\
GPT-6 Luna    & 0.619 [0.539, 0.696] & \textbf{0.771} & 0.802 & 24 & 11 & \$0.22 \\
\bottomrule
\end{tabular}
\end{center}\vspace{-6pt}
\begin{itemize}
\item \textbf{No significant accuracy differences.}
All fit within each other's confidence interval.
\item \textbf{Different pathologies}. GLM almost never abstains (highest accuracy,
  lowest precision). Haiku and Luna abstain more and are right more often when they commit.
\item \textbf{Safety filters cost Luna the most.} OpenAI's bio-risk filter blocked 15
  questions at least once and Luna declined 8 (mostly virus-related); Haiku was blocked on
  1. On the 87 questions with no block or refusal, all three are within 3 pts (GLM 0.697,
  Luna 0.687, Haiku 0.667).
\item \textbf{There is little spread within a single model / single question.} Five repeats
  narrowed CIs only from about $\pm$9 to $\pm$8 pts. Caveats: public split only; possible
  training contamination.
\end{itemize}

\sect{RL environments (BioProBench): proposal}
\textbf{Tasks.} We propose turning four single-turn BioProBench tasks (Liu et al., 2026;
CC BY-NC) into environments, since each has gold labels that can be checked automatically.
Training uses the Hugging Face train splits, minus 1,000 validation items per task. The
official test splits are held out (overlap with train is $\le$3 items per task). Protocol
generation (GEN) is excluded because its only rewards are reference-similarity scores (BLEU,
embeddings), which are easy to game.

\vspace{-2pt}\begin{center}\small
\begin{tabular}{@{}lp{4.3cm}rp{6.4cm}@{}}
\toprule
Task & Model sees $\rightarrow$ returns & Train & Proposed reward (in $[0,1]$) \\
\midrule
PQA & Protocol sentence with a blank, 5 options $\rightarrow$ option + confidence & 47,083 &
  $0.9\cdot\mathrm{correct} + 0.1\cdot(1-(c/100-\mathrm{correct})^2)$: correctness plus
  a small calibration term \\
ERR & One step + purpose and neighboring steps $\rightarrow$ True/False & 93,049 &
  1 if the verdict matches the label, else 0 \\
ORD & Shuffled steps $\rightarrow$ index permutation & 27,549 &
  $0.5\cdot\mathrm{exact} + 0.5\cdot(\tau+1)/2$ (Kendall's $\tau$; the order of identical
  steps is ignored) \\
REA & ERR + written analysis of operation, reagent, and parameter errors & (ERR) &
  $\mathrm{correct}\cdot(0.7+0.3\cdot\mathrm{judge})$ on items that contain an error;
  the LLM-judge bonus never rescues a wrong verdict \\
\bottomrule
\end{tabular}
\end{center}\vspace{-6pt}
\begin{itemize}
\item \textbf{Training signal.} All rewards score the final answer only, with no
  credit for intermediate steps. PQA and ERR are essentially binary. ORD gives partial
  credit. REA is the only signal on reasoning, and it is noisier because it relies on a
  judge. Only the last answer block counts, and missing, malformed, or truncated answers
  score 0, so the reward also enforces the answer format. Parsing is strict:
  ERR accepts only an exact \texttt{True}/\texttt{False}, because the original scorer counts any
  answer containing ``true'' as True (copying the prompt's ``True or False'' template
  scores $\approx$50\%).
\item \textbf{Data cleaning.} We would drop 3,228 ERR train items (7\% of the error items)
  that are labeled erroneous but show the corrected step, 200 step texts that carry both
  labels, malformed items, and duplicates. This should leave about 46k PQA, 87k ERR, and
  26k ORD items.
\item \textbf{Reward tests (before training).} (i) The gold answer for every item should
  score 1. (ii) Unit tests should give 0 to known exploits: template echo, ``untrue'',
  1-based or partial index lists, missing or repeated tags. (iii) Random and constant
  policies should land at the no-skill floors (PQA 0.2, ERR 0.5, ORD at the computed random
  level). (iv) The rewards should agree with the original BioProBench scorers on
  well-formed outputs.
\item \textbf{Training runs.} GRPO on a small open instruct model (1--4B parameters). Runs:
  a zero-shot baseline, one single-task run each for PQA, ERR, and ORD, a mixed-task run,
  and REA with and without the judge bonus. Each run stops early on mean validation
  reward (higher is better, $\Delta_{\min}=0.005$, patience 5 evaluations), and we keep the
  best checkpoint. During training we would watch per-class ERR accuracy (to catch
  collapse to one label), parse failures, output length, KL to the base model, and the
  gap between the judge score and verdict accuracy for REA (a sign of judge gaming).
\item \textbf{Evaluation.} BioProBench test results with the paper's metrics, plus
  \emph{the ProtocolQA eval above}, to test whether BioProBench training transfers to
  troubleshooting on a separate benchmark.
\end{itemize}

\sect{Next steps (one more week)}
\begin{enumerate}[leftmargin=1.4em]
\item \textbf{Separate filtering from ability.} Re-ask blocked and refused items through
  other endpoints, then report both ``as deployed'' and ``answered-only'' scores.
\item \textbf{Open-answer scoring.} Two off-the-shelf judges map free-text answers to the
  stored options, validated on known answers, to test whether the options do the work.
\item \textbf{Effort sweep and error taxonomy.} Run \texttt{low}/\texttt{high}/\texttt{max}
  effort to get accuracy-vs-cost curves. Tag items by error type (omission, quantity,
  ordering, condition) to see where each model fails.
\item \textbf{More power.} Add related public tasks; 108 items cannot resolve $\sim$6-pt gaps.
\end{enumerate}
\end{document}
